\documentclass[11pt,a4paper]{article}
\usepackage[T1]{fontenc}
\usepackage[utf8]{inputenc}
\usepackage[brazilian]{babel}
\usepackage{lmodern}
\usepackage[margin=2.3cm]{geometry}
\usepackage{enumitem}
\usepackage{amssymb}

\setlength{\parindent}{0pt}
\setlength{\parskip}{0.55em}
\pagestyle{plain}
\newcommand{\caixa}{$\square$\ }
\newcommand{\campo}[1]{\rule{#1}{0.4pt}}

\begin{document}

\begin{center}
  {\large\bfseries Termo de Consentimento Livre e Esclarecido (TCLE)}\\[0.3em]
  {\small Entrevista com gestores de estabelecimentos e fornecedores}
\end{center}

\textbf{Título da pesquisa:} Cardápios digitais com autoatendimento em bares e restaurantes de Porto Alegre: um estudo de caso sobre adoção, motivações e desafios pós-2020.

\textbf{Pesquisador responsável:} César Hoffmann, aluno do curso de Bacharelado em Ciência da Computação da Universidade do Vale do Rio dos Sinos (UNISINOS). Contato: cesarhoff@unisinos.br, +55 51 9 9804-4896 (telefone e WhatsApp).

\textbf{Orientador:} Prof.\ Dr.\ Guilherme Lacerda, Escola Politécnica, UNISINOS.

\textbf{1. Objetivo.} Analisar como bares e restaurantes de Porto Alegre percebem os benefícios, riscos e desafios técnicos e operacionais do uso, ou da escolha de não usar, cardápios digitais com autoatendimento desde 2020, e como seus clientes vivenciam essa interação.

\textbf{2. Procedimento.} Entrevista semiestruturada de cerca de 45 a 60 minutos, presencial ou por videoconferência, gravada em áudio ou vídeo para posterior transcrição.

\textbf{3. Riscos e benefícios.} Os riscos são mínimos e limitam-se a eventual desconforto ao falar de aspectos operacionais do negócio. Não há benefício financeiro. O benefício esperado é contribuir para o conhecimento sobre transformação digital no setor de alimentação e para a melhoria das soluções disponíveis.

\textbf{4. Garantias.} A participação é voluntária e pode ser interrompida a qualquer momento, sem justificativa e sem prejuízo. O participante pode pedir que trechos da entrevista não sejam usados e tem acesso aos resultados consolidados mediante solicitação.

\textbf{5. Dados pessoais (LGPD, Lei 13.709/2018).} Os dados são usados apenas para fins acadêmicos ligados ao Trabalho de Conclusão de Curso e a eventuais publicações dele decorrentes. Gravações e transcrições ficam em ambiente de acesso restrito ao pesquisador e ao orientador e são descartadas em até 24 meses após a defesa. O participante pode pedir, a qualquer momento, acesso, correção ou exclusão de seus dados. O nome do participante não é divulgado no trabalho.

\textbf{6. Identificação no trabalho.} Marque uma opção:

\caixa Autorizo que o nome do estabelecimento ou da empresa apareça no trabalho.\\
\caixa Prefiro que o estabelecimento ou a empresa seja identificado apenas por um pseudônimo (por exemplo, ``Caso D'').

\textbf{7. Declaração de consentimento.} Declaro que li e compreendi as informações acima, que tive a oportunidade de esclarecer dúvidas com o pesquisador e que consinto livremente em participar desta pesquisa e em ter a entrevista gravada nos termos descritos.

\vspace{1.2em}
Nome do participante: \campo{10.5cm}

Estabelecimento ou empresa: \campo{9.6cm}\qquad

Data da entrevista: \campo{4cm}

\vspace{2.5em}
\campo{7.2cm}\hfill\campo{7.2cm}\\
Assinatura do participante\hfill Assinatura do pesquisador

\vfill
{\small Este termo pode ser assinado à mão e digitalizado, ou assinado digitalmente (por exemplo, pela assinatura eletrônica gratuita do gov.br). Uma via fica com o participante e outra com o pesquisador.}

\end{document}
